\section{Introduction}

A city road network is not static. Signals, incidents, construction and
closures change the cost of travel continuously, and a routing method that is
optimal on a frozen snapshot may be worthless if it has to recompute from
nothing every time the network changes. This observation motivates a large
literature on \emph{bio-inspired} routing, in which ant colony optimization
(ACO)~\cite{dorigo1996ant,stutzle2000maxmin} and models of the slime mould
\emph{Physarum polycephalum}~\cite{nakagaki2000maze,tero2007mathematical}
maintain an internal state (a pheromone field, a field of tube
conductivities) that is claimed to let them \emph{adapt} to a change
rather than restart.

The claim is attractive and it is rarely measured. Evaluations of these
methods overwhelmingly use synthetic graphs, report accuracy on an
undisrupted shortest-path problem, and treat the adaptive machinery as
self-evidently useful. Three gaps follow. Test graphs are typically
undirected and dense, so they lack the structure (one-way streets, low out-degree, cul-de-sacs)
that makes real road routing hard. Congestion is
usually absent, so the edge weights that the adaptation is supposed to track
never move. And recovery after a disruption, the very behaviour that
justifies the adaptive state, is seldom the measured quantity.

This paper closes those gaps for one city. We benchmark classical Dijkstra and
A*~\cite{dijkstra1959note,hart1968formal}, a MAX--MIN Ant System (MMAS), and
a Tero-style Physarum conductivity model on the real road network of central Bengaluru under congestion generated by
the SUMO microsimulator~\cite{lopez2018microscopic}, and we measure recovery
from seven classes of road closure under both warm start (retained internal
state) and cold start (state reset). The design is fully paired: every
algorithm sees identical origin--destination pairs, identical closures and
identical traffic, and all quality figures are expressed against a Dijkstra
optimum computed on the same weighted graph.

\subsection{Contributions}

\textbf{1) Path extraction, not the model, decides what a Physarum result
means.} The conductivity model does not emit a route; a separate extraction
step does. We instrument that step and find that Physarum's apparent
competence is largely attributable to its fallback. Pooled over all
extraction methods it matches the optimal route on 90.0\% of undisrupted
pairs, reproducing the accuracy reported in the literature; restricted to its
own maximum-flux walk the figure is 23.3\%. The median stretch is exactly
1.000 inside \emph{every} provenance class, so a pooled number cannot
distinguish ``the model found the optimum'' from ``the fallback returned
Dijkstra's answer''. On warm recoveries the effect is starker still: the
conductivity threshold subgraph never connects origin to destination, and the
model reports plain Dijkstra's path on 66.8\% of cells. We therefore report
Physarum as two labelled series throughout and never pool them.

\textbf{2) Warm starting buys speed and charges for it in quality.}
Retaining the conductivity field across a closure reduces recovery by
2.97$\times$ in iterations and 6.87$\times$ in wall-clock, paired on
identical cells ($n=358$, $p=3.2\times10^{-59}$ and
$9.0\times10^{-52}$). The distribution is bimodal rather than merely
shifted: 48.3\% of warm recoveries converge in exactly one iteration, which cold start never achieves; its minimum is 73. Per-iteration cost is
unchanged (0.051\,s versus 0.054\,s), which is what establishes that the
saving is adaptation and not cheaper arithmetic. Against this, warm start is
significantly \emph{worse} on the reported post-disruption path
($p=4.3\times10^{-7}$). Reporting only the speed-up, as the self-healing
framing invites, would describe half of the effect.

\textbf{3) Adaptation improves extraction quality while reducing its
reliability.} Against a zero-adaptation control (the same pressure solve on a uniform
conductivity field, containing none of the Tero dynamics), the converged
field completes a walk \emph{less} often (43.3\% versus 56.7\%)
and yet produces far better routes when it does (median stretch 1.006 versus
1.437; exact match 46.2\% versus 5.9\% of completions). Completion and
quality move in opposite directions, so either measured alone misleads. This
also explains the discrepancy with prior theory: a near-uniform flux field is
easy to walk precisely because it is uninformative.

\textbf{4) The sign of the congestion effect belongs to the congestion
model.} Under an analytic BPR volume-delay function congestion \emph{helps}
ACO; under SUMO microsimulation on the same network it \emph{hurts} it
(within-10\% rate 81.3\% light versus 58.0\% heavy, paired $p=0.036$). Any
claim about behaviour ``under congestion'' that does not name its congestion
model is underdetermined.

Throughout we report what the methods lose as well as what they gain.
Dijkstra recomputes an exact route in a single iteration and roughly 31\,ms,
and no bio-inspired method here approaches that; the defensible quantities
are iterations-to-adapt and post-disruption path quality, which is the
speed-versus-quality trade this paper measures.
